\section{这门课的形式：Executable Lecture}

课程里有一个很有特色的说法：这是一门 executable lecture。
意思不是“讲义里有代码”这么简单，而是：讲义本身就像一个可执行的程序，内容、代码和结构是统一的。

这有两个象征性的好处：
\begin{itemize}
    \item 你可以把学习当作“运行一个系统”，而不是只读静态文档。
    \item 课程本身就体现了“构建即理解”的方法论。
\end{itemize}

\begin{lstlisting}[language=Python,caption={课程式讲义的概念性示意}]
理解系统  ->  实现系统  ->  运行系统  ->  观察系统  ->  继续改进
\end{lstlisting}

